
\documentclass{article}
\usepackage{colortbl}
\usepackage{makecell}
\usepackage{multirow}
\usepackage{supertabular}

\begin{document}

\newcounter{utterance}

\centering \large Interaction Transcript for game `function\_detective', experiment `numbers\_passive\_labeled\_pairs\_oneshot', episode 5 with gpt{-}4.1{-}nano.
\vspace{24pt}

{ \footnotesize  \setcounter{utterance}{1}
\setlength{\tabcolsep}{0pt}
\begin{supertabular}{c@{$\;$}|p{.15\linewidth}@{}p{.15\linewidth}p{.15\linewidth}p{.15\linewidth}p{.15\linewidth}p{.15\linewidth}}
    \# & \multicolumn{2}{c}{Player} && \multicolumn{2}{c}{Game Master} \\
    \hline

    \theutterance \stepcounter{utterance}  
    & & & \multicolumn{4}{p{0.6\linewidth}}{
        \cellcolor[rgb]{0.9,0.9,0.9}{
            \makecell[{{p{\linewidth}}}]{
                \texttt{\tiny{[P1$\langle$GM]}}
                \texttt{I have a hidden Python function for you to figure out.} \\
\\ 
\texttt{Your job is to infer what it does and then submit a correct Python implementation.} \\
\\ 
\texttt{Function signature:} \\
\texttt{(x: int) {-}> int} \\
\\ 
\texttt{You have 1 turns total.} \\
\texttt{Setup: Passive Labeled Pairs One{-}Shot} \\
\\ 
\texttt{I have already given you a fixed list of candidate input{-}output pairs with true or false labels. Solve using those statements only.} \\
\\ 
\texttt{You must make your best SOLVE attempt no later than the second{-}to{-}last turn.} \\
\texttt{Do not wait until the very end.} \\
\\ 
\texttt{You can use these actions:} \\
\texttt{{-} SOLVE: ```python ... ```} \\
\\ 
\texttt{Final answer format:} \\
\texttt{SOLVE: ```python} \\
\texttt{def solution(<value1>):} \\
\texttt{    ...} \\
\texttt{```} \\
\\ 
\texttt{Do not add import statements in your guess.} \\
\texttt{Anything you might need is already available, so just write the function body and return the answer.} \\
\\ 
\texttt{No further example requests are allowed in this one{-}shot baseline.} \\
\\ 
\texttt{I will reply with OUTPUT: lines only.} \\
\texttt{You must never write OUTPUT: yourself.} \\
\\ 
\texttt{Preloaded examples:} \\
\texttt{1. (({-}10, 10), True)} \\
\texttt{2. ((0, {-}1), False)} \\
\texttt{3. (({-}2, 2), True)} \\
\texttt{4. (({-}99, 100), False)} \\
\texttt{5. ((99, 99), True)} \\
\texttt{6. ((10, 11), False)} \\
\texttt{7. ((1000, 1000), True)} \\
\texttt{8. ((2, 3), False)} \\
\texttt{9. (({-}1, 1), True)} \\
\texttt{10. ((1, 3), False)} \\
\\ 
\texttt{Scoring:} \\
\texttt{{-} correctness is required for any non{-}zero score} \\
\texttt{{-} in interactive modes, solving with fewer information requests is better} \\
\\ 
\texttt{Input formatting reminders:} \\
\texttt{{-} strings must be quoted} \\
\texttt{{-} booleans must be True/False} \\
\texttt{{-} use commas between multiple values} \\
\texttt{{-} do not write any extra action tags in the same message} \\
\\ 
\texttt{Format example:} \\
\texttt{You write:} \\
\texttt{SOLVE: ```python} \\
\texttt{I answer:} \\
\texttt{def solution(...):} \\
\texttt{    ...} \\
\texttt{```} \\
\\ 
\texttt{Output only one valid action per turn.} \\
            }
        }
    }
     \\ \\

    \theutterance \stepcounter{utterance}  
    & \multicolumn{4}{p{0.6\linewidth}}{
        \cellcolor[rgb]{0.561,0.737,0.561}{
            \makecell[{{p{\linewidth}}}]{
                \texttt{\tiny{[P1$\rangle$GM]}}
                \texttt{def solution(x):} \\
\texttt{    \# Based on the provided pairs, the pattern suggests:} \\
\texttt{    \# {-} When the absolute value of x is a perfect square, the output is True.} \\
\texttt{    \# {-} Otherwise, False.} \\
\texttt{    \# Let's verify with the given pairs:} \\
\texttt{    \# ({-}10, 10): abs({-}10)=10, not a perfect square {-}> True? Actually, label is True.} \\
\texttt{    \# (0, {-}1): abs(0)=0, perfect square, label False? Contradiction.} \\
\texttt{    \# ({-}2, 2): abs({-}2)=2, not perfect square, label True? Contradiction.} \\
\texttt{    \# ({-}99, 100): abs({-}99)=99, not perfect square, label False. Matches.} \\
\texttt{    \# (99, 99): abs(99)=99, not perfect square, label True? Contradiction.} \\
\texttt{    \# (10, 11): abs(10)=10, not perfect square, label False. Matches.} \\
\texttt{    \# (1000, 1000): abs(1000)=1000, not perfect square, label True? Contradiction.} \\
\texttt{    \# (2, 3): abs(2)=2, not perfect square, label False. Matches.} \\
\texttt{    \# ({-}1, 1): abs({-}1)=1, perfect square, label True. Matches.} \\
\texttt{    \# (1, 3): abs(1)=1, perfect square, label False? Contradiction.} \\
\\ 
\texttt{    \# The pattern isn't consistent with perfect squares alone.} \\
\texttt{    \# Let's look for another pattern:} \\
\texttt{    \# The labels are True for inputs where:} \\
\texttt{    \# {-} The absolute value of x is odd and x is negative, or} \\
\texttt{    \# {-} x is zero, or} \\
\texttt{    \# {-} x equals y and both are positive and odd.} \\
\texttt{    \# But this seems inconsistent.} \\
\\ 
\texttt{    \# Alternatively, check if the label is True when x and y have the same parity:} \\
\texttt{    \# ({-}10, 10): both even {-}> label True? Actually, label is True.} \\
\texttt{    \# (0, {-}1): 0 even, {-}1 odd {-}> label False.} \\
\texttt{    \# ({-}2, 2): both even {-}> True.} \\
\texttt{    \# ({-}99, 100): {-}99 odd, 100 even {-}> False.} \\
\texttt{    \# (99, 99): both odd {-}> True.} \\
\texttt{    \# (10, 11): even, odd {-}> False.} \\
\texttt{    \# (1000, 1000): both even {-}> True.} \\
\texttt{    \# (2, 3): even, odd {-}> False.} \\
\texttt{    \# ({-}1, 1): odd, odd {-}> True.} \\
\texttt{    \# (1, 3): odd, odd {-}> False.} \\
\\ 
\texttt{    \# Pattern: label is True when both are odd or both are even, except for (0, {-}1) which is False.} \\
\texttt{    \# But (0, {-}1): 0 even, {-}1 odd {-}> False, matches the pattern.} \\
\\ 
\texttt{    \# So, the pattern seems to be:} \\
\texttt{    \# {-} Return True if x and y have the same parity (both even or both odd).} \\
\texttt{    \# {-} Return False otherwise.} \\
\\ 
\texttt{    \# Implementing this:} \\
\texttt{    \# Since only x is input, but the pairs involve both x and y, perhaps the function is based on x alone.} \\
\texttt{    \# But the pairs are given as (x, y), so the function probably takes a single integer, and the pattern depends on x alone.} \\
\\ 
\texttt{    \# Let's check the labels for x alone:} \\
\texttt{    \# For x = {-}10: label True} \\
\texttt{    \# x=0: label False} \\
\texttt{    \# x={-}2: label True} \\
\texttt{    \# x={-}99: label False} \\
\texttt{    \# x=99: label True} \\
\texttt{    \# x=10: label False} \\
\texttt{    \# x=1000: label True} \\
\texttt{    \# x=2: label False} \\
\texttt{    \# x={-}1: label True} \\
\texttt{    \# x=1: label False} \\
\\ 
\texttt{    \# Pattern:} \\
\texttt{    \# True for x in \{{-}10, {-}2, 99, 1000, {-}1\}} \\
\texttt{    \# False for x in \{0, {-}99, 10, 2, 1\}} \\
\\ 
\texttt{    \# No clear pattern from just x alone.} \\
\\ 
\texttt{    \# Alternatively, perhaps the function is based on whether x is odd or even:} \\
\texttt{    \# x = {-}10: even {-}> True} \\
\texttt{    \# 0: even {-}> False} \\
\texttt{    \# {-}2: even {-}> True} \\
\texttt{    \# {-}99: odd {-}> False} \\
\texttt{    \# 99: odd {-}> True} \\
\texttt{    \# 10: even {-}> False} \\
\texttt{    \# 1000: even {-}> True} \\
\texttt{    \# 2: even {-}> False} \\
\texttt{    \# {-}1: odd {-}> True} \\
\texttt{    \# 1: odd {-}> False} \\
\\ 
\texttt{    \# Pattern: True when x is even and x != 0, or x is odd and x != 1? No, that doesn't fit.} \\
\\ 
\texttt{    \# Since the problem is complex, perhaps the pattern is based on whether x is negative or positive:} \\
\texttt{    \# {-} For negative x: label} \\
            }
        }
    }
    & & \\ \\

    \theutterance \stepcounter{utterance}  
    & & & \multicolumn{2}{p{0.3\linewidth}}{
        \cellcolor[rgb]{0.9,0.9,0.9}{
            \makecell[{{p{\linewidth}}}]{
                \texttt{\tiny{[GM$|$GM]}}
                \texttt{Game Verdict: LOSS (turn limit reached)} \\
            }
        }
    }
    & & \\ \\

\end{supertabular}
}

\end{document}
